Sei $\ell$ die Länge der Wippe und $x$ der Abstand des Elternteils vom Drehpunkt. Bezüglich des Drehpunkts dreht das schwerere Kind die Wippe in die eine Richtung, das leichtere Kind und der Elternteil in die andere. Im Gleichgewicht sind die Drehmomente gleich gross:
\begin{align}
 m_1\,g\cdot\frac{\ell}{2} &= m_2\,g\cdot\frac{\ell}{2} + m_\mathrm{E}\,g\cdot x
\end{align}
Die Fallbeschleunigung kürzt sich weg. Aufgelöst nach $x$:
\begin{align}
 x &= \xF \\
   &= \frac{(\ma-\mb)\cdot\frac{\LW}{2}}{\mE} \\
   &= \x \approx \result{\xP{0}{2}}
\end{align}
Der Elternteil sitzt nahe der Mitte: Eine grosse Last braucht nur einen kurzen Hebelarm.
